\documentclass[a4paper, 12pt]{article}
\usepackage[english, russian]{babel}
\usepackage[T2A]{fontenc}
\usepackage[utf8]{inputenc}
\usepackage{indentfirst}
\usepackage{geometry}
\usepackage{titleps}
\usepackage{mathtext}
\usepackage{amsmath}
\usepackage{enumitem}
\usepackage{tabularx}
\usepackage{tikz}
\usetikzlibrary{positioning}
\setlength{\tabcolsep}{4pt}
\newpagestyle{main}{
\setheadrule{0pt}
\sethead{}{}{}
\setfootrule{0pt}
\setfoot{}{}{}
}
\pagestyle{main}

\usepackage{amsthm}
\geometry{top=2cm}
\begin{document}
\begin{center}
{\large КОНСПЕКТ ЛЕКЦИЙ}

ПО ДИСЦИПЛИНЕ <<ТЕОРЕТИЧЕСКАЯ ИНФОРМАТИКА>>

\vspace{5pt}
\textbf{Основы теории информации}
\end{center}
\vspace{10pt}

\section*{Информация и данные}

\textbf{Информация} --- это совокупность сведений, полученных путём наблюдения
и изучения, расширяющих представление об объектах, явлениях окружающей среды,
их свойствах, состоянии и взаимосвязях.

\textbf{Данные} --- это информация в форме, пригодной для её хранения,
обработки, передачи.

Основоположник теории информации --- \underline{Шеннон}.

\section*{Система передачи информации}

Обмен информацией происходит между людьми, между людьми и окружающим миром.
Информацию передаёт какой-либо объект, система.

\begin{center}
\resizebox{\textwidth}{!}{%
\begin{tikzpicture}[
  node distance=2.2cm,
  box/.style={draw, minimum height=1.1cm, minimum width=2cm, align=center},
  lbl/.style={font=\footnotesize, align=center},
  >=stealth]
  \node[box] (src) {Источник\\информации};
  \node[box, right=of src] (tx) {Передатчик};
  \node[box, right=of tx] (ch) {Канал\\связи};
  \node[box, right=of ch] (rx) {Приёмник};
  \node[box, right=of rx] (dst) {Потребитель\\информации};
  \node[box, below=1.3cm of ch] (noise) {Источник\\помех};
  \draw[->] (src) -- node[lbl, above] {сообщение} (tx);
  \draw[->] (tx) -- node[lbl, above] {передаваемый\\сигнал} (ch);
  \draw[->] (ch) -- node[lbl, above] {принятый\\сигнал} (rx);
  \draw[->] (rx) -- node[lbl, above] {сообщение} (dst);
  \draw[->] (noise) -- node[lbl, right] {помехи} (ch);
\end{tikzpicture}}
\end{center}

\textbf{Источник} отправляет информацию в виде сообщения.

\textbf{Сообщение} --- это форма представления информации для её
непосредственной передачи в виде данных. Формы сообщений:
% ниже список в оригинале читается лишь частично; восстановлен по смыслу
\begin{itemize}[noitemsep]
  \item текстовая --- текст на естественном языке;
  \item кодовая --- закодированные данные;
  \item графическая --- изображения объектов;
  \item звуковая;
  \item видео;
  \item сигналы.
\end{itemize}

\textbf{Передатчик} преобразует сообщение в сигналы, пригодные для передачи
по каналу связи. Например, в мобильной телефонии передатчик преобразует
звуковое сообщение в радиосигналы.

\textbf{Сигнал} (в общем случае) --- это физический динамический процесс,
при котором его параметры изменяются во времени.

\textbf{Канал связи} --- это физическая среда, используемая для передачи
сигнала от передатчика к приёмнику. Каналом может быть коаксиальный кабель,
полоса радиочастот, луч света и т.\,д.

\textbf{Приёмник} обычно выполняет преобразование, обратное преобразованию,
проводимому передатчиком. Восстанавливает сообщение из сигнала.

\textbf{Потребитель информации} --- это человек или устройство, для которого
предназначено сообщение.

Чаще всего при передаче сигнала эти сигналы могут быть искажены помехами.
\underline{Источник помех} действует на передаваемые сигналы, изменяя
принимаемое сообщение.

Система передачи информации проявляется в различных сферах нашей жизни:

\begin{center}
\begin{tabularx}{\textwidth}{|X|X|X|X|}
\hline
\textbf{Части} & \textbf{Лекция} & \textbf{Запрос веб-сервера (ВС)} & \textbf{Телевидение} \\
\hline
Источник информации & Лектор & Клиентский компьютер (КК) & Источник радиосигнала \\
\hline
Передатчик & Голосовые связки лектора & Сетевая плата КК & Радиоканал \\
\hline
Канал связи & Воздушная среда & Провайдер, канал связи (ПКС) & Воздух, радиоканал \\
\hline
Приёмник & Уши & Сетевая плата & Антенна радиоприёмника \\
\hline
Потребитель информации & Студент & ВС & Слушатель \\
\hline
Источник помех & Шум в аудитории, уличный шум & Воздействие на ПКС & Атмосферные шумы \\
\hline
Сообщение & Текст & Текстовое или числовое & Видеоизображение \\
\hline
\end{tabularx}
\end{center}

\section*{Дискретные и непрерывные сигнальные каналы}

Если параметр сигнала принимает конечное число значений и при этом они все
могут быть пронумерованы, то этот сигнал называется \underline{дискретным}.
Канал для передачи таких сигналов также называется дискретным.

Примером дискретных сигналов являются буквы. Количество букв конечно, и их
можно рассматривать как уровни сигнала передачи сообщения.

Если параметр сигнала является непрерывной во времени функцией, то сигнал
называется \underline{непрерывным}.

\textbf{Пример:} человеческая речь --- передаваемая звуковой волной,
характеризуемой частотой, фазой и амплитудой. Параметром сигнала в этом
случае является давление, создаваемое этой волной в точке нахождения
приёмника --- человеческого уха. Канал для передачи непрерывных сигналов
называется непрерывным.

Для передачи информации используются непрерывные и дискретные каналы связи.
Сигналы переводят из непрерывных в дискретные и обратно.

\section*{Дискретизация}

% Страница с дискретизацией читается хуже остальных; сомнительные места помечены.
\textbf{Дискретизация} --- это процесс преобразования непрерывного сигнала
в дискретный.

Для этого диапазон значений функции (ось ординат) разбивается на конечное
число отрезков равной ширины, называемой \underline{шагом дискретизации}.
Дискретное значение --- новое значение отрезка, в который попало непрерывное
значение функции. % формулировка восстановлена по смыслу

Чем меньше шаг дискретизации, тем точнее дискретный сигнал соответствует
непрерывному сигналу. % последняя строка в оригинале неразборчива

\end{document}
